\documentclass[10pt,a4paper]{amsart}
\usepackage[margin=19mm]{geometry}
\usepackage{amsmath,amssymb,mathtools}
\usepackage[scr=rsfs,cal=euler]{mathalfa}
\usepackage{xfrac}
\usepackage[unicode,hidelinks,bookmarks=false]{hyperref}
\newcommand{\ie}{i.e.\ }
\newcommand{\cf}{cf.\ }
\newcommand{\resp}{respectively\ }
\newcommand{\Subsection}{\subsection}
\providecommand{\nobreakdash}{\nobreak\hbox{-}}
\setlength{\emergencystretch}{2em}
\multlinegap0pt
\usepackage{amssymb}
\usepackage{bm}
\usepackage{breqn}
\usepackage{float}
\usepackage{tikz}
\usepackage{xcolor}
\usepackage{enumerate}
\newtheorem{proposition}{Proposition}
\title{Zeros of the generalized Wronskian--Hermite polynomials}
\author{Chengfa Wu, Guangxiong Zhang}
\date{Source: arXiv:2607.29027v1 (CC BY 4.0)}
\begin{document}
\maketitle

\noindent\emph{Source and licence.} Zeros of the generalized Wronskian--Hermite polynomials. Version arXiv:2607.29027v1 (CC BY 4.0), distributed by arXiv under the Creative Commons Attribution 4.0 International licence, http://creativecommons.org/licenses/by/4.0/, as recorded in the arXiv licence metadata for this exact version and shown on the abstract page https://arxiv.org/abs/2607.29027v1. Full licence notice: \texttt{licenses/arxiv-2607.29027-CC-BY-4.0.txt}.\par\smallskip

\noindent\textit{Editorial scope.} Counted unit: the article's proposition prop:k4-W-recurrence (source0.tex lines 1789--1802). The article's complete original proof of it is delivered with it (source0.tex lines 1805--2089, one \texttt{proof} environment of 285 lines). No mathematics is written, repaired or re-derived: the statement and the proof are the article's own lines, in the article's own notation, and no retained line is substituted or re-flowed.  Retained uncounted as the source-local context the proof needs: lines 114--115.  Nothing was omitted from any retained span, so the package carries no omission record.  No retained line contains a citation, so the wrapper carries no bibliography and no bibliography entry.  No retained line contains a figure environment, an image-inclusion command or drawing code, so no image is delivered and the package declares no asset dependency.  Editorial formatting only, in addition: an AMS \texttt{amsart} preamble that restates the article's own declarations and macros used by the retained lines, namely the environments \texttt{proposition} on separate counters, together with the standard packages those lines need.  The retained environments print in this document as Proposition 1 in order of appearance, whereas the article prints the same items with the numbering of its own full document.  The complete native source of the article as distributed in the arXiv e-print for this exact version is bundled unchanged as \texttt{sources/206455/source0.tex}.
\begin{equation}\label{eq:p-generating-intro} \sum_{n=0}^{\infty}p_n^{[m]}(z)\epsilon^n = \exp(z\epsilon+\epsilon^m).
\end{equation} 

\begin{proposition}
\label{prop:k4-W-recurrence}
For $N\geq1$ and $l \in \{1,2,3\}$, the generalized Wronskian--Hermite polynomials \(W_N^{[2,4,l]}(z)\) satisfy the recurrence relation
\begin{eqnarray}
\label{eq:k4-W-recurrence}
    W_{N+1}^{[2,4,l]}(z)W_{N-1}^{[2,4,l]}(z)&=&8z\left(W_N^{[2,4,l]}(z)\left(W_N^{[2,4,l]}\right)''(z)-\left(\left(W_N^{[2,4,l]}\right)'(z)\right)^2\right) \nonumber\\
    &+&8W_N^{[2,4,l]}(z)\left(W_N^{[2,4,l]}\right)'(z)+z\left(z^2+12N+4l-6\right)\left(W_N^{[2,4,l]}(z)\right)^2,
\end{eqnarray}
with initial polynomials
\begin{equation}
\label{eq:k4-W-initial-values}
    W_0^{[2,4,l]}(z)=1,\qquad W_1^{[2,4,1]}(z)=z,\qquad W_1^{[2,4,2]}(z)=z^2+2,\qquad W_1^{[2,4,3]}(z)=z^3+6z.
\end{equation}
\end{proposition}

\begin{proof}
Let \(l\in\{1,2,3\}\) and \(N\geq1\) be fixed. Denote 
\begin{equation}
    \Delta_N^{(l)}(z)=\operatorname{Wr} \left[p_{n_1}^{[2]}(z),p_{n_2}^{[2]}(z),\ldots,p_{n_{N-1}}^{[2]}(z),p_{n_N}^{[2]}(z)\right],\;
    \Sigma_N^{(l)}(z)=\operatorname{Wr} \left[p_{n_1}^{[2]}(z),p_{n_2}^{[2]}(z),\ldots,p_{n_{N-1}}^{[2]}(z),p_{n_{N+1}}^{[2]}(z)\right],
\end{equation}
where $n_j=l+4(j-1),\; j=1,2,\ldots,N+1$, and $$ \operatorname{Wr} [f_1,f_2, \dots, f_N] = \det \left( f_j^{(i-1)}(z) \right)_{i,j=1}^N.$$ By the Jacobi formula for determinants, we obtain
\begin{equation}\label{eqn:Jacobi formula for determinants}
    \Delta_{N+1}^{(l)}(z)\Delta_{N-1}^{(l)}(z)
    =
    \Delta_N^{(l)}(z)\left(\Sigma_N^{(l)}\right)'(z)
    -
    \left(\Delta_N^{(l)}\right)'(z)\Sigma_N^{(l)}(z).
\end{equation}
Furthermore, recalling the relation $W_N^{[2,4,l]}(z)=c_N^{[2,4,l]}\Delta_N^{(l)}(z)$ and multiplying both sides by \(c_{N+1}^{[2,4,l]}c_{N-1}^{[2,4,l]}\), we obtain
\begin{equation}
\label{eq:k4-normalized-Jacobi-corrected}
    W_{N+1}^{[2,4,l]}(z)W_{N-1}^{[2,4,l]}(z)
    =
    W_N^{[2,4,l]}(z)\left(\widetilde{\Sigma}_N^{(l)}\right)'(z)
    -
    \left(W_N^{[2,4,l]}\right)'(z)\widetilde{\Sigma}_N^{(l)}(z),
\end{equation}
where
\begin{equation}
    \widetilde{\Sigma}_N^{(l)}(z)
    =
    \frac{c_{N+1}^{[2,4,l]}c_{N-1}^{[2,4,l]}}{c_N^{[2,4,l]}}\Sigma_N^{(l)}(z) = \frac{(n_N+1)(n_N+2)(n_N+3)(n_N+4)}{4N} c_N^{[2,4,l]}\Sigma_N^{(l)}(z).
\end{equation}


From the definition of $p_n^{[2]}(z)$ from \eqref{eq:p-generating-intro}, its differential and recurrence relations are governed by
\begin{equation}
    \left(p_n^{[2]}\right)'(z)=p_{n-1}^{[2]}(z),
    \qquad
    (n+1)p_{n+1}^{[2]}(z)=zp_n^{[2]}(z)+2\left(p_n^{[2]}\right)'(z).
\end{equation}
Consequently, we deduce that
\begin{equation}
    z\left(p_n^{[2]}\right)'(z)=n p_n^{[2]}(z)-2p_{n-2}^{[2]}(z),
\end{equation}
and
\begin{equation}
    z^2p_n^{[2]}(z)=(n+1)(n+2)p_{n+2}^{[2]}(z)-(4n+2)p_n^{[2]}(z)+4p_{n-2}^{[2]}(z).
\end{equation}
Applying the latter identity repeatedly yields
\begin{equation}
\begin{aligned}
    z^4p_n^{[2]}(z)
    ={}&(n+1)(n+2)(n+3)(n+4)p_{n+4}^{[2]}(z)-4(n+1)(n+2)(2n+3)p_{n+2}^{[2]}(z) \\
    &+12(2n^2+2n+1)p_n^{[2]}(z)
      -16(2n-1)p_{n-2}^{[2]}(z)+16p_{n-4}^{[2]}(z).
\end{aligned}
\end{equation}
Combining the last two formulas, we get
\begin{equation}
\label{eq:k4-one-step-column-identity}
    \left(32z\frac{d}{dz}+z^4+(8n+12)z^2+8n^2+8n+12\right)p_n^{[2]}(z)
    =(n+1)(n+2)(n+3)(n+4)p_{n+4}^{[2]}(z)+16p_{n-4}^{[2]}(z).
\end{equation}

We now define the auxiliary sum of determinants $\mathcal{S}_N^{(l)}(z)$ as
\begin{equation}
\label{eq:k4-SN-definition}
    \mathcal S_N^{(l)}(z)=\sum_{j=1}^{N}\operatorname{Wr} \left[p_{n_1}^{[2]},\ldots,p_{n_{j-1}}^{[2]},\left(32z\frac{d}{dz}+z^4+(8n_j+12)z^2+8n_j^2+8n_j+12\right)p_{n_j}^{[2]},p_{n_{j+1}}^{[2]},\ldots,p_{n_N}^{[2]}\right].
\end{equation}
On the one hand, substituting \eqref{eq:k4-one-step-column-identity} with $n=n_j$, the replaced column in the $j$-th Wronskian takes the form
\begin{equation}
    (n_j+1)(n_j+2)(n_j+3)(n_j+4)p_{n_j+4}^{[2]}(z)+16p_{n_j-4}^{[2]}(z).
\end{equation}
For \(1\leq j\leq N-1\), we have $p_{n_j+4}^{[2]}=p_{n_{j+1}}^{[2]}$, which implies that the first term creates two identical columns, causing the corresponding Wronskian to vanish. Similarly, for \(2\leq j\leq N\), the relation $p_{n_j-4}^{[2]}=p_{n_{j-1}}^{[2]}$
ensures that the second term also yields two identical columns. Finally, for $j=1$, we have
\begin{equation}
    p_{n_1-4}^{[2]}=p_{l-4}^{[2]}=0,
    \qquad
    l=1,2,3.
\end{equation}
As a result, only the term  \(p_{n_N+4}^{[2]}=p_{n_{N+1}}^{[2]}\) in the last summand survives, which gives
\begin{equation}
\label{eq:k4-SN-to-Sigma}
    \mathcal S_N^{(l)}(z)
    =
    (n_N+1)(n_N+2)(n_N+3)(n_N+4)\Sigma_N^{(l)}(z).
\end{equation}

On the other hand, we can express \(\mathcal S_N^{(l)}(z)\) in terms of \(\Delta_N^{(l)}(z)\) and its derivative. 
To facilitate this calculation, we denote
\begin{equation}
    B_j(z)=z^4+(8n_j+12)z^2+8n_j^2+8n_j+12.
\end{equation}
Applying the Leibniz rule, for  \(r=0,1,\ldots,N-1\), the $(r+1)$-th row of the \(j\)-th replaced column in \(\mathcal S_N^{(l)}\) can be expanded as
\begin{equation}\label{eqn:expantion of the j-th column}
\begin{aligned}
    \frac{d^r}{dz^r}\left[32z\left(p_{n_j}^{[2]}\right)' +B_j p_{n_j}^{[2]} \right] &=
    32z\left(p_{n_j}^{[2]}\right)^{(r+1)} 
    +32r\left(p_{n_j}^{[2]}\right)^{(r)} 
    +B_j \left(p_{n_j}^{[2]}\right)^{(r)} +rB_j' \left(p_{n_j}^{[2]}\right)^{(r-1)} 
     \\
    &+\binom{r}{2}B_j'' \left(p_{n_j}^{[2]}\right)^{(r-2)} 
    +\binom{r}{3}B_j^{(3)} \left(p_{n_j}^{[2]}\right)^{(r-3)} 
    +\binom{r}{4}B_j^{(4)} \left(p_{n_j}^{[2]}\right)^{(r-4)},
\end{aligned}
\end{equation}
where, for \(r<q\), the term containing 
\(\left(p_{n_j}^{[2]}\right)^{(r-q)}(z)\) is understood to be zero, and the derivatives of $B_j(z)$ are given by
\begin{equation}
    B_j'(z)=4z^3+(16n_j+24)z,\quad
    B_j''(z)=12z^2+16n_j+24,\quad
    B_j^{(3)}(z)=24z,\quad B_j^{(4)}(z)=24.
\end{equation}

For the subsequent determinant manipulations, given a column vector \(\mathbf{a}=(a_0,\ldots,a_{N-1})^T\), we denote by
\(\Delta_N^{(l)}[j;\mathbf{a}]\) the determinant obtained from \(\Delta_N^{(l)}\) by replacing its \(j\)-th column with \(\mathbf{a}\). Differentiating the Wronskian \(\Delta_N^{(l)}(z)\) yields 
\begin{equation}
\label{eq:k4-column-derivative-rule}
    \left(\Delta_N^{(l)}\right)'(z)=\sum_{j=1}^{N}
    \Delta_N^{(l)}\!\left[j;\left(\left(p_{n_j}^{[2]}\right)^{(r+1)}(z)\right)_{r=0}^{N-1}\right],
\end{equation}
where we use the convention
\begin{equation}
    \left(\alpha\left(r\right)\right)_{r=0}^{N-1}= \left(\alpha(0),\alpha(1),\cdots,\alpha(N-1)\right)^T.
\end{equation}
If the column vector \(\mathbf{a}\) does not depend on \(j\), then we have
\begin{equation}
\label{eq:k4-diagonal-row-rule}
    \sum_{j=1}^{N}
    \Delta_N^{(l)}\!\left[j;\left(a_r\left(p_{n_j}^{[2]}\right)^{(r)}(z)\right)_{r=0}^{N-1}\right]
    =
    \left(\sum_{r=0}^{N-1}a_r\right)\Delta_N^{(l)}(z).
\end{equation}
Similarly, if the column vector \(\mathbf{a}\) does not depend on \(j\), for \(s\geq1\), we get
\begin{equation}
\label{eq:k4-lower-row-rule}
    \sum_{j=1}^{N}
    \Delta_N^{(l)}\!\left[j;\left(a_r\left(p_{n_j}^{[2]}\right)^{(r-s)}(z)\right)_{r=0}^{N-1}\right]
    =
    0,
\end{equation}
where the corresponding element of the
replacement column is understood to be $0$ for $r<s$. Moreover, we also have
\begin{equation}
\label{eq:k4-r-upper-row-rule}
    \sum_{j=1}^{N}
    \Delta_N^{(l)}\!\left[j;\left(r\left(p_{n_j}^{[2]}\right)^{(r+1)}(z)\right)_{r=0}^{N-1}\right]
    =
    (N-1)\left(\Delta_N^{(l)}\right)'(z).
\end{equation}
Furthermore, from the relation
\begin{equation}
    2\left(p_n^{[2]}\right)''(z)+z\left(p_n^{[2]}\right)'(z)-np_n^{[2]}(z)=0,
\end{equation}
we get
\begin{equation}
\label{eq:k4-differentiated-Hermite-corrected}
    2\left(p_n^{[2]}\right)^{(s+2)}(z)
    +z\left(p_n^{[2]}\right)^{(s+1)}(z)
    +(s-n)\left(p_n^{[2]}\right)^{(s)}(z)
    =
    0.
\end{equation}
Taking \(n=n_j\) and \(s=r-1\), for \(r\geq1\), gives
\begin{equation}
    r n_j\left(p_{n_j}^{[2]}\right)^{(r-1)}(z)
    =
    2r\left(p_{n_j}^{[2]}\right)^{(r+1)}(z)
    +rz\left(p_{n_j}^{[2]}\right)^{(r)}(z)
    +r(r-1)\left(p_{n_j}^{[2]}\right)^{(r-1)}(z).
\end{equation}
Using \eqref{eq:k4-diagonal-row-rule}-\eqref{eq:k4-r-upper-row-rule}, we obtain
\begin{equation}
\label{eq:k4-nj-first-lower-shift}
\sum_{j=1}^{N}
    \Delta_N^{(l)}\!\left[j;\left(rn_j\left(p_{n_j}^{[2]}\right)^{(r-1)}(z)\right)_{r=0}^{N-1}\right]
    =
    2(N-1)\left(\Delta_N^{(l)}\right)'(z)
    +\frac{N(N-1)}{2}z\Delta_N^{(l)}(z).
\end{equation}
Similarly, taking \(n=n_j\) and \(s=r-2\), for \(r\geq2\), we have
\begin{equation}
    n_j\left(p_{n_j}^{[2]}\right)^{(r-2)}(z)
    =
    2\left(p_{n_j}^{[2]}\right)^{(r)}(z)
    +z\left(p_{n_j}^{[2]}\right)^{(r-1)}(z)
    +(r-2)\left(p_{n_j}^{[2]}\right)^{(r-2)}(z),
\end{equation}
which implies that
\begin{equation}
\label{eq:k4-nj-second-lower-shift}
\begin{aligned}
     \sum_{j=1}^{N}
    \Delta_N^{(l)}\!\left[j;\left(\binom{r}{2}n_j\left(p_{n_j}^{[2]}\right)^{(r-2)}(z)\right)_{r=0}^{N-1}\right]=
    2\sum_{r=0}^{N-1}\binom{r}{2}\Delta_N^{(l)}(z)
    =
    \frac{N(N-1)(N-2)}{3}\Delta_N^{(l)}(z).
\end{aligned}
\end{equation}
Thus, we can collect all contributions in the expansion of \(\mathcal S_N^{(l)}(z)\) by using \eqref{eqn:expantion of the j-th column}, which results in
\begin{equation}
\label{eq:k4-SN-collected-before-sums}
    \mathcal S_N^{(l)}(z)
    =
    32Nz\left(\Delta_N^{(l)}\right)'(z)
    +
    \Bigg[
        \sum_{j=1}^{N}B_j(z)
        +8N(N-1)z^2
        +16N(N-1)
        +\frac{16N(N-1)(N-2)}{3}
    \Bigg]\Delta_N^{(l)}(z),
\end{equation}
where
\begin{equation}
\begin{aligned}
    \sum_{j=1}^{N}B_j(z)
    =
    Nz^4
    +\left(8\sum_{j=1}^{N}n_j+12N\right)z^2
    +8\sum_{j=1}^{N}n_j^2
    +8\sum_{j=1}^{N}n_j
    +12N.
\end{aligned}
\end{equation}
Since
\begin{equation}
    \sum_{j=1}^{N}n_j
    =
    \sum_{j=1}^{N}\left(l+4(j-1)\right)
    =
    N(l+2N-2),
\end{equation}
and
\begin{equation}
    \sum_{j=1}^{N}n_j^2
    =
    \sum_{j=1}^{N}\left(l+4(j-1)\right)^2=
    \frac{N}{3}\left(16N^2+(12l-24)N+3l^2-12l+8\right),
\end{equation}
substituting these sums into \eqref{eq:k4-SN-collected-before-sums}, we obtain
\begin{equation}
\label{eq:k4-SN-final}
\begin{aligned}
    \mathcal S_N^{(l)}(z)
    =
    32Nz\left(\Delta_N^{(l)}\right)'(z)
    +
    \Big[
        Nz^4
        +4N(6N+2l-3)z^2
        +4N\left(12N^2+(8l-12)N+2l^2-6l+3\right)
    \Big]\Delta_N^{(l)}(z).
\end{aligned}
\end{equation}

Finally, combining \eqref{eq:k4-SN-to-Sigma} and \eqref{eq:k4-SN-final}, we have
\begin{equation}
\label{eq:k4-Sigma-tilde-corrected}
    \widetilde{\Sigma}_N^{(l)}(z)
    =
    8z\left(W_N^{[2,4,l]}\right)'(z)
    +
    \left[
        \frac{z^4}{4}
        +(6N+2l-3)z^2
        +12N^2+(8l-12)N+2l^2-6l+3
    \right]W_N^{[2,4,l]}(z),
\end{equation}
and
\begin{equation}\label{eq:k4-Sigma-tilde-Prime-corrected}
\begin{aligned}
    \left(\widetilde{\Sigma}_N^{(l)}\right)'(z)
    =
    8\left(W_N^{[2,4,l]}\right)'(z)
    +8z\left(W_N^{[2,4,l]}\right)''(z)
    +\left(z^3+2(6N+2l-3)z\right)W_N^{[2,4,l]}(z) \\
    +
    \left[
        \frac{z^4}{4}
        +(6N+2l-3)z^2
        +12N^2+(8l-12)N+2l^2-6l+3
    \right]\left(W_N^{[2,4,l]}\right)'(z).
\end{aligned}
\end{equation}
Substituting \eqref{eq:k4-Sigma-tilde-corrected} and \eqref{eq:k4-Sigma-tilde-Prime-corrected} into \eqref{eq:k4-normalized-Jacobi-corrected}, we obtain \eqref{eq:k4-W-recurrence}.
This completes the proof.
\end{proof}

\end{document}
